\chapter{لوحات المعلومات والتقارير}
\label{ch:builder-dashboards}

\section{ماذا ستبني في هذا الفصل؟}
\label{sec:builder-dashboards-intro}

بعد أن تتراكم السجلات يحتاج المديرون إلى أرقام تلخّص العمل، مثل عدد الطلبات المفتوحة أو المتأخرة. ويشرح هذا الفصل بناء لوحة المعلومات من أولها إلى عرضها: تنشئها في القائمة، ثم تضيف إليها الأقسام والعناصر وتربط كل عنصر بمصدر بياناته، وتضبط إعداداتها ومرشحاتها، ثم تعاينها وتنشرها. ويشرح بعد ذلك كيف يعرضها الفريق ويصدّرها، وكيف تُنشر لوحات الذكاء الاصطناعي صفحاتٍ في التطبيقات، وكيف تطرح أسئلة سريعة في تقارير الذكاء الاصطناعي.

ويرى كل مشاهد في اللوحة ما تسمح به صلاحياته ما لم يختر صاحبها غير ذلك صراحة، لذلك تتكرر في الصفحات مراجعة مستوى الوصول قبل النشر.

\begin{rcnote}[مثال: شركة الأفق]
يبني مسؤول النظام في الأفق لوحة «متابعة الزيارات» تعرض طلبات الزيارة في كل حالة وفي كل فرع، ثم ينشرها في تطبيق «خدمات الفروع»، فتراجع سارة الخطيب كل صباح الطلبات التي تنتظر قرارها.
\end{rcnote}

\section{صفحات هذا الفصل}
\label{sec:builder-dashboards-pages}

تبدأ أربع صفحات من قسم «الأساسية» في أدوات البناء: «لوحات المعلومات» و«الصفحات» و«لوحات معلومات الذكاء الاصطناعي» و«تقارير الذكاء الاصطناعي». أما الصفحات الثماني بينها فتشرح منشئ اللوحة الذي تفتحه من بطاقة اللوحة، ثم عرض اللوحة المنشورة. واتبع الترتيب في المرة الأولى، لأن كل صفحة تبني على ما قبلها.

\begin{table}[H]
\centering
\caption{التقارير والمتابعة}
\label{tab:builder-pages-reports}
\rcZebra
\begin{tabularx}{\linewidth}{R{4.2cm} Y}
\rcHeaderRow
\thh{الصفحة وقسمها} & \thh{ماذا تقدم هذه الصفحة؟}\\
لوحات المعلومات — الأساسية & تنشئ اللوحات وترتّبها في مجلدات وتحدد من يصل إلى كل منها.\\
منشئ اللوحة: الأقسام والعناصر — منشئ اللوحة & تضيف الأقسام والعناصر وترتّبها في الشبكة وتحفظ المسودة.\\
إعدادات العنصر ومصادر بياناته — منشئ اللوحة & تختار مصدر أرقام العنصر وطريقة حسابها ومستوى الوصول ومظهره.\\
أنواع العناصر — منشئ اللوحة & تعرّف بالعناصر والرسوم البيانية المتاحة ومتى تستخدم كلاً منها.\\
إعدادات اللوحة — منشئ اللوحة & تحدد اسم اللوحة ومكانها في التنقل والبيانات التي يراها المشاهد.\\
مرشحات المشاهد ومحددات الكيانات — منشئ اللوحة & تضيف مرشحات التاريخ والسجلات والخيارات التي يغيّرها المشاهد.\\
منشئ اللوحة بالذكاء الاصطناعي — منشئ اللوحة & تصف اللوحة بكلماتك فيقترح الذكاء الاصطناعي عناصرها لتراجعها.\\
المعاينة والنشر والإصدارات — منشئ اللوحة & تعاين اللوحة وتنشرها وتعود إلى إصدار سابق عند الحاجة.\\
عرض اللوحة وتصديرها — عرض اللوحة & يغيّر المشاهد المرشحات ويصدّر بيانات جدول أو تقرير \en{PDF}.\\
الصفحات — الأساسية & تنشر لوحات الذكاء الاصطناعي صفحاتٍ مستقلة داخل التطبيقات.\\
لوحات معلومات الذكاء الاصطناعي — الأساسية & تبني لوحات معلومات بمساعدة الذكاء الاصطناعي في مساحة خاصة بك.\\
تقارير الذكاء الاصطناعي — الأساسية & تجيب عن أسئلتك عن بيانات المؤسسة بتقرير مكتوب.\\
\end{tabularx}
\end{table}
